\documentclass[UTF8,11pt]{ctexart}
\usepackage{amsmath,amssymb,mathtools,geometry}
\geometry{a4paper,margin=20mm}\allowdisplaybreaks
\begin{document}
\section*{按紧急优先级加权的受灾点物资分配模型}

\paragraph{来源与边界} PDF 第10--12页，§3.4，式(17)--(25)。固定一个决策周期 $t$；
需求预测、区域分组、紧急度和上期库存均是本次优化的已知输入。下一阶段供应调拨不属于本模型。
\subsection*{集合、参数和变量}
$L$ 为物资种类，$M$ 为配送中心集合，$G$ 为受灾区域组，$I_g$ 为组 $g$ 的区域集合，
$V$ 为运输车辆类型。各组 $I_g$ 不交。
$D_i^l(t)\ge0$ 为需求，$Q_m^l(t-1|t-1)\ge0$ 为中心期初库存；
$UNV^l$ 为物资单位体积，$UNL_v$ 为车型容量，$NV_v^m(t)$ 为本期可用车次。
$CS_m^l$、$CT_{mi}^l$ 分别为单位设置和运输成本；$u_i^1(t)$ 为距上次物资到达的时间。
$\Omega_g(t)$ 为输入的组紧急评价值（旧转写为 $X_g(t)$，此处按原 PDF 使用 $\Omega_g(t)$）。
已知归一化界满足 $\overline F_g^j>\underline F_g^j$，$j=1,2$；
$\omega_r,\omega_c\ge0$ 且 $\omega_r+\omega_c=1$。
只对总需求 $D_g=\sum_{l\in L,i\in I_g}D_i^l(t)>0$ 的有效组优化，零需求组的出货由需求约束固定为零；
要求 $\sum_g\Omega_g(t)\ne0$。紧急权重和目标分量定义为
\begin{align}
\gamma_g(t)&=\frac{\Omega_g(t)}{\sum_{h\in G}\Omega_h(t)},\\
F_g^1(X)&=\frac{\sum_{l\in L,m\in M,i\in I_g}X_{mi}^l(t)}{D_g},\\
F_g^2(X)&=\sum_{l\in L,m\in M,i\in I_g}
  \bigl[CS_m^l u_i^1(t)+CT_{mi}^l\bigr]X_{mi}^l(t),\\
F_g(X)&=\omega_r\frac{F_g^1(X)-\underline F_g^1}{\overline F_g^1-\underline F_g^1}
-\omega_c\frac{F_g^2(X)-\underline F_g^2}{\overline F_g^2-\underline F_g^2}.
\end{align}
唯一决策族 $X_{mi}^l(t)\ge0$ 是中心 $m$ 向区域 $i$ 分配的物资 $l$ 数量。
\subsection*{完整优化模型}
\begin{align}
\max_X\quad &\sum_{g\in G}\gamma_g(t)F_g(X)\tag{20}\\
\text{s.t.}\quad
&\sum_{m\in M}X_{mi}^l(t)\le D_i^l(t)
 &&g\in G,\ i\in I_g,\ l\in L,\tag{22}\\
&\sum_{g\in G,i\in I_g}X_{mi}^l(t)\le Q_m^l(t-1|t-1)
 &&m\in M,\ l\in L,\tag{23}\\
&\sum_{g\in G,i\in I_g,l\in L}UNV^lX_{mi}^l(t)
\le\sum_{v\in V}UNL_vNV_v^m(t) &&m\in M,\tag{24}\\
&X_{mi}^l(t)\ge0 &&m\in M,\ g\in G,\ i\in I_g,\ l\in L.\tag{25}
\end{align}
固定 $t$ 后这是线性规划。输出供后续供应调拨及库存更新使用，但后续目标不反向改变本次分配。
旧转写式(24)将车辆类型与时段都写为 $t$，此处将车辆类型统一写为 $v$；原 PDF 使用不同字形的车辆类型索引。

\end{document}
